

% ---------------
We reuse the expensive quantities to reduce computation.
$(\bX^T \bV \bX)^{-1}$, $\bP\by$ and $\bC\bP\by$ are repeatedly used in equations (\ref{reml_gradient}) and (\ref{avg_information}).
After getting (\ref{gls_fe}), we get the the residual $\by - \bX\bbeta^{(n)}$.
$\bP\by$ follows from the identity $\bP^{(n)}\by = \bV^{(n)-1} (\by - \bX \bbeta^{(n)})$,
and have already computed $(\bX^T \bV^{(n)} \bX)^{-1}$ from the usual matrix inversion.
$\bC \bP \by$ follows from a single matrix-vector product that we explain shortly.

The efficient genetic relatedness matrix-vector multiplication implemented in \textsf{tskit} is the key ingredient of our algorithm.
Let $\bC = \bZ \Sigma_\bu \bZ^T$ be the individual genetic relatedness matrix.
\textsf{GeneticRelatednessMatrix} implements $\bv \mapsto \bC\bv$.
A more commonly used subroutine is its extension 
$\textsf{CovarianceMatrix}(\sigma, \tau, \bv) = (\sigma^2 \bI + \tau^2 \bC ) \bv$ which trivially follows from \textsf{GeneticRelatednessMatrix}.
We can multiply the inverse of $\sigma^2 \bI + \tau^2 \bC $ by conjugate gradient deploying \textsf{CovarianceMatrix}.
We denote this as $\textsf{CovarianceMatrix}^{-1}(\sigma, \tau, \bv) = (\sigma^2 \bI + \tau^2 \bC)^{-1} \bv$.

\begin{taocpalg}{D}{First and second derivatives} \label{algo:D}
{
    $\by$ is the trait vector and $\bX$ is the fixed effects covariates.
    Compute the gradients (\ref{reml_gradient}) and the average informations (\ref{avg_information}).
    $\tau^{(n)}$, $\sigma^{(n)}$ and $\bbeta^{(n)}$ are the working variance components and the fixed effects coefficient at step $n$.
    Define $\bV^{(n)}$ and $\bP^{(n)}$ in the same fashion.
}


\algstep{D1.}{Compute $\bbeta^{(n)}$}{ 
    Compute equation (\ref{gls_fe}). 
    Keep $(\bX^T \bV^{(n)} \bX)^{-1}$,  
    $\bP^{(n)}\by$, and $\bC \bP^{(n)}\by$ for uses in later steps.
j}

\algstep{D2.}{Compute gradient}{\\
    Compute $\Tr(\bP^{(n)})$ and  
    $\bP^{(n)} \by = \textsf{CovarianceMatrix}^{-1}(\sigma^{(n)}, \tau^{(n)}, \by - \bX \bbeta^{(n)})$.\\
    $\bC^{(n)}\bP^{(n)} \by = \textsf{CovarianceMatrix}(0, 1, \bP^{(n)}\by)$. \\
    $\by^T \bP^{(n)} \bP^{(n)} \by = (\bP^{(n)} \by)^T(\bP^{(n)} \by)$. \\
    $\by^T \bP^{(n)} \bC^{(n)} \bP^{(n)} \by = (\bC^{(n)}\bP^{(n)} \by)^T(\bP^{(n)} \by)$.
}

\algstep{D3.}{Compute average information}{ \\
    $\bV^{-1(n)} \bP^{(n)}\by = \textsf{CovarianceMatrix}^{-1}(\sigma^{(n)}, \tau^{(n)}, \bP^{(n)}\by)$. \\
    $\bV^{-1(n)} \bC^{(n)} \bP^{(n)}\by = \textsf{CovarianceMatrix}^{-1}(\sigma^{(n)}, \tau^{(n)}, \bC^{(n)}\bP^{(n)}\by)$. \\
    %
    $\by^T \bP^{(n)} \bV^{-1(n)} \bP^{(n)}\by = 
    (\bV^{-1(n)}\bP^{(n)} \by)^T \bP^{(n)}\by
    $. \\
    $\by^T \bP^{(n)} \bC^{(n)} \bV^{-1(n)} \bC^{(n)} \bP^{(n)}\by = 
    (\bV^{-1(n)} \bC^{(n)} \bP^{(n)} \by)^T \bC^{(n)}\bP^{(n)}\by.
    $ \\
    %
    $\bX^T \bV^{-1(n)} \bP^{(n)}\by = \bX^T  \cdot \bV^{-1(n)} \bP^{(n)}\by$. \\
    $\bX^T\bV^{-1(n)} \bC^{(n)} \bP^{(n)}\by = \bX^T \cdot \bV^{-1(n)} \bC^{(n)} \bP^{(n)}\by$.\\
    %
    $\by^T \bP^{(n)} \bV^{-1(n)} \bX (\bX^T \bV^{-1(n)}\bX)^{-1} \bX^T \bV^{-1(n)} \bP^{(n)}\by \\
    = \by^T \bP^{(n)} \bV^{-1(n)} \bX \cdot (\bX^T \bV^{-1(n)}\bX)^{-1} \cdot \bX^T \bV^{-1(n)} \bP^{(n)}\by$ . \\
    $\by^T \bP^{(n)} \bC^{(n)} \bV^{-1(n)} \bX (\bX^T \bV^{-1(n)}\bX)^{-1} \bX^T \bV^{-1(n)}  \bP^{(n)}\by \\
    = \by^T \bP^{(n)} \bC^{(n)} \bV^{-1(n)} \bX \cdot (\bX^T \bV^{-1(n)}\bX)^{-1} \cdot \bX^T \bV^{-1(n)} \bP^{(n)}\by$.  \\
    $\by^T \bP^{(n)} \bC^{(n)} \bV^{-1(n)} \bX (\bX^T \bV^{-1(n)}\bX)^{-1} \bX^T \bV^{-1(n)} \bC^{(n)} \bP^{(n)}\by \\
    = \by^T \bP^{(n)} \bC^{(n)} \bV^{-1(n)} \bX \cdot (\bX^T \bV^{-1(n)}\bX)^{-1} \cdot \bX^T \bV^{-1(n)} \bC^{(n)} \bP^{(n)}\by$. \\
}

\end{taocpalg}

\textsf{CovarianceMatrix}$^{-1}$ dominates most of the computation in \textbf{Algorithm D}.
This routine occurs twice in gradient computation and twice more in average information computation.
Therefore, the overall routine is only as twice as expensive compared to computing gradient only.
This additional cost is payed off (supplementary figures?)


The concerns around confounding and false discoveries caused by population structure in genome-wide association studies led to the development of
numerous methods such as genomic control, principal component (PC) adjustment, and linear mixed models \citep{devlin1999genomiccontrol, price2006eigenstrat, kang2010variancecomponent}.
Principal component adjustment and linear mixed models continue to be popular in literature,
and recently proposed methods tend to include the elements of both approaches \citep{loh2015boltlmm, jiang2019fastgwa}.

Early discussions centered on determining the most suitable method for adjusting population structure \citep{price2010popstrat, sul2013selection}.
Shortly after, people started to recognize the close relationship between the two methods, 
as both methods used the same genetic relatedness matrix \citep{hoffman2013lmm}.
Linear mixed model includes the entire relatedness matrix,
and PC adjustment takes the top eigenvectors of the relatedness matrix after eigendecomposition.
For this reason,
some authors argued that PC adjustment approximates the linear mixed model \citep{hoffman2013lmm, yao2023limitationpca}.
A recent analysis shows that this claim is only partially correct
because the linear mixed model and PC adjustment weigh the eigenvectors
differently \citep{schraiber2024unifying}.
These analyses are based on the eigenspectrum of the genetic relatedness matrix.

An alternative viewpoint comes from the interpretation of genetic effects as either fixed or random effects.
It is common to test a single SNP at a time,
and other SNPs that may be potentially associated with the trait
are omitted from the regression \citep{sul2018confounding}.
If an omitted causal variant is correlated with the tested variant,
the tested variant appears to be associated with the trait even in the absence of true causation \citep{vilhjlmsson2012natureconfounding, sul2018confounding, veller2024interpreting}.
Thus, one must account for the omitted variants in univariate regression.
PC adjustment can be thought of as including the omitted variants in the regression by reducing their dimension via principal component analysis.
PC covariates produced from dimension reduction are included as fixed effects in the regression.
Liner mixed models integrate over the omitted variable by treating them as random variables \citep{sul2018confounding}.

A common feature of these different viewpoints on the nature of confounding in genome-wide association studies
is that they are either purely mathematical (eigenspectrum) or statistical (fixed versus random effects),
where the arguments are only weakly connected to the underlying evolutionary process.
Some efforts have been made to clarify that population structure itself is not the problem, 
but the allele frequency stratification caused by genetic drift is the real problem in association testing \citep{vilhjlmsson2012natureconfounding, edge2013windfalls}.

Here, we present a unifying view that combines all the aforementioned perspectives from an evolutionary perspective based on ARG-LMM.
As we've stressed earlier, 
genetic effects are in fact branch effects and not the variants' effect.
Branch effects are the aggregate effect of mutations that have occurred on it,
and the randomness is solely driven by the random mutational process.
Similar to how this interpretation led to evolutionary interpretations of random effects,
we can explain how evolutionary forces are related to confounding, in particular the fixed effects, in complex trait analysis.

Surprisingly, the fixed components of the genetic effects $\bZ\bfix$ turn out to be absent when the evolution is neutral.
Under neutral evolutionary scenarios that may or may not accompany demographic events that cause population structure, 
we may assume that the mutation rate is constant for each position \citep{wakeley2008coalescent, durrett2008probabilitydna},
i.e., $u_{ep} = u_p$ for all $e: p\in e$ (edge $e$ contains position $p$). 
In this case, we show that $\bZ\bfix$ reduces to a constant.
\begin{align} 
    \begin{aligned}
        \left[\bZ \bfix \right]_n 
        &= \sum_{e=1}^E \bZ_{ne} \E{ \sum_{p:p\in e} \bbeta_p \mathbf{1}_{ep}}  
        = \sum_{e=1}^E \sum_{p:p\in e} \bZ_{ne} \bbeta_p l_eu_p  \\
        &= \sum_{p=1}^P \sum_{e:p\in e} \bZ_{ne} \bbeta_p l_eu_p 
        = \sum_{p=1}^P \bbeta_p u_p \left(\sum_{e:p\in e} \bZ_{ne}  l_e \right) \\ 
        &= \sum_{p=1}^P \bbeta_p u_p \cdot 2t_{\mathrm{root},p}
        = \text{const. resp. to $n$}
    \end{aligned}
\end{align}
This means that the fixed component of the genetic effect is fully accounted for by an intercept.
Hence, once we include an intercept in the regression model, 
further adjustment by PC covariates is unnecessary. 
This argument formalizes earlier claims that population structure itself is not a concern \citep{vilhjlmsson2012natureconfounding, edge2013windfalls}.
Regardless of how complicated the demographic events and the population structure are, 
their directional contribution in the regression as fixed effects is vacant.
Nonetheless, genetic drift does not play a role here, as claimed earlier, because conditioning on the genealogy wipes the variation of drift out of the inference. 
In sum, fixed effects covariates are obsolete for correcting genetic confounding if the evolutionary process is neutral.

The previous result does not, however, preclude the utility of PC adjustment in genetic association studies for two reasons.
First of all, rarely is an evolutionary process purely neutral.
Various studies have shown that although selective sweeps are rare in the human genome,
negative and background selection is universal [cite].
This means that we can no longer assume a constant mutation rate at a given position across edges.
This will cause the fixed portion of genetic effects $\bZ\bfix$ to be non-trivial, unlike the neutral scenario.
This is in line with evidence from simulations and empirical studies where pure linear mixed models without PC covariates
are not properly calibrated in populations affected by selection \citep{price2010popstrat, sul2013selection, price2013response, vilhjlmsson2012natureconfounding, edge2013windfalls}.


The time zero of evolutionary change must be specified in any evolutionary analysis.
For example, the notion of identity-by-descent (IBD) is valid given a particular reference time point because any genetic materials will eventually coalesce if we go sufficiently further back in time \citep{Thompson2013}.
The idea of fixing a time point in which the evolutionary process began has a long history that dates back to \cite{Wright1949} and has been expanded in subsequent studies in the context of quantitative genetics \citep{weir1990gda, weir2002fstatistics, Ochoa2021, hou2023kinshiprobust}.
Choosing this time point is also important in ARG-LMM's generative model.
The current paper makes an implicit choice of which we condition at
the time of the oldest root node in the tree sequence, i.e., the grand TMRCA.
However, one could choose another time point below the grand TMRCA.
All mutations that occurred above this time point but below the grand TMRCA will be conditioned and become fixed effects.
These effects are not constant across individuals even if the evolution is neutral.
If this time is set below the grand TMRCA, 
one should include all the genetic effects above this time 
as fixed effects in the regression model.
The relatedness driven by random effects will decay as we get closer to the present and eventually vanish as soon as we arrive.
This reiterates the previous arguments that relatedness is a relative measure with respect to a reference time point that is set by the researcher,
and no single measure of relatedness exists in an absolute sense \citep{Thompson2013}.
We plan to add an option in \textsf{tslmm} to set this time point and automatically handle the evolution of relatedness and fixed effects as a function of time zero.

We can further generalize equation (\ref{origin_fe})'s argument to incorporate time-varying mutation rate within an edge.
In this setting, we stipulate that a mutation rate at a position only depends on time, i.e., the mutation rate is a function of time $u_p(t)$.
\begin{align} 
    \begin{aligned}
        \left[\bZ \bfix \right]_n 
        &= \sum_{e=1}^E \bZ_{ne} \E{ \sum_{p:p\in e} \bbeta_p \mathbf{1}_{ep}}  
        = \sum_{e=1}^E \sum_{p:p\in e} \bZ_{ne} \bbeta_p \E{ \mathbf{1}_{ep}}  \\
        &= \sum_{e=1}^E \sum_{p:p\in e} \bZ_{ne} \bbeta_p \int_{t(e(c))}^{t(e(p))} u_p(t) dt  
        = \sum_{p=1}^P \bbeta_p \sum_{e: p \in e } \bZ_{ne}  \int_{t(e(c))}^{t(e(p))} u_p(t) dt  \\
        &=  2 \cdot \sum_{p=1}^P \bbeta_p \int_{0}^{t_{\mathrm{root},p}} u_p(t) dt
        = \text{const. resp. to $n$}
    \end{aligned}
\end{align}
Therefore, a non-trivial fixed effect is present only if there is a heterogeneity in the mutation rate across both time and lineages.
As a side note, our argument only discusses the role of fixed effects as a correction for genetic confounding.
PC covariates may correct non-genetic confounders that align with the population structure, and so would be useful despite not being relevant to correcting genetic confounding.





Unlike the usual definition of additive variance, this quantity does not depend on the frequency of alleles.
Indeed, the relationship between $\tau^2$ and additive genetic variance is convoluted and we discuss this relationship in depth.

Recall the close resemblance between the CoA model and ARG-LMM.
A core quantity of the CoA model is the mutational variance 
\begin{align}
\sigma_m^2 = \sum_l \sum_{p:p \in l} \bbeta_p^2 u_l
\label{eq:mutational_variance}
\end{align}
where $\sigma_l^2$ and $u_l$ are the allelic effect's variance and the mutation rate of locus $l$, respectively \citep{kimura1964numberofalleles,lynch1998quant, walsh2018evolution}.
\cite{lynch1986phenoneutral} have shown that in a neutrally evolving additive trait of a diploid population of size $N_e$,
the additive genetic variance coincides with $4N_e\sigma_m^2$  under the drift-mutation equilibrium only in \textit{expectation}.
When genetic drift is non-negligible, as much as the frequency of individual variants, the additive variance of the trait fluctuates randomly over time and across evolutionary replicates.
On the other hand, the effective population size $N_e$ and the mutational variance $\sigma_m^2$ are fixed parameters that govern the evolution of the system.
Hence, the additive genetic variance can only be equal to a combination of fixed numbers $4N_e \sigma_m^2$ in \textit{average}.
This means that rarely is a population's additive variance a fixed parameter in an evolutionary system incorporating genetic drift.
Note that increasing the population size nor the sample size remove this variability \citep{lynch1986phenoneutral}.
The same issue is present in ARG-LMM which describes the evolution until the time of full coalescence.
% avg. coalesence time ~Ne, drift rate 1-1/Ne => converge to constant factor ~e^-1

To cope with the randomness of additive genetic variance, one could simulate a trait with a fixed additive variance $V_G$  by the following steps.
After simulating the population's genetic value $\bg$, we divide the genetic values by its population standard deviation $\sqrt{\var{\bg}} = \sqrt{\sum_{n=1}^N (\bg_n - \overline{\bg})^2 / (N-1)}$.
Then, multiplying the genetic values by $\sqrt{V_G}$ will give the desired result.
To further specify the narrow sense heritability $h^2$, adding a normal distribution with variance $(1-h^2)/h^2 \cdot V_G$ will do. 
This procedure has been pursued earlier in \cite{zhang2021biobank} and \cite{tagami2024tstrait}.
However, a careful inspection of the above process tells that it is impossible to recover $V_G$ nor $h^2$ in practice.
Suppose that a trait was generated from an ARG-LMM with $\tau^2$ and $\sigma^2$. 
Scaling by $V_G / \var{\bg}$ and $(1-h^2)/h^2 \cdot V_G$ of the aforementioned procedure turns the ARG-LMM to have new parameters $\tau^{2}_{\textrm{new}} = \tau^2 \cdot V_G / \var{\bg}$ and $\sigma^{2}_{\textrm{new}} = \sigma^2 \cdot (1-h^2)/h^2 \cdot V_G$.
An unbiased estimator (e.g. \texttt{tslmm}) recovers only the new parameters $\tau^2_{\textrm{new}}$ and $\sigma^2_{\textrm{new}}$ which are functions of three unknown quantities $V_G$, $h^2$, and $\var{\bg}$.
That said, we cannot recover any of the three quantities from data as we only have two equations.
The benchmark of estimating $h^2$ presented in \cite{zhang2021biobank} was only able to recover $V_G$ and $h^2$ due to knowing the value of $\bg$ in the simulation.
Hence, the benchmark does not inform real world analysis that has no access to the true genetic values $\bg$.

We demonstrate this argument through neutral evolution experiments repeated multiple times, mimicking the theoretical setup of \cite{lynch1986phenoneutral}.
We have $\sigma_m^2 = \tau^2 \cdot P$ due to the resemblance between equations (\ref{eq:global_tau}) and (\ref{eq:mutational_variance}) (see \textbf{Materials and Methods}).
Thus, we can estimate the mutational variance using \texttt{tslmm} through multiplying the total sequence length to the $\tau^2$ estimate.
As shown in \textbf{Figure \ref{fig:mutationalvariance}A}, $P \widehat{\tau}_{\texttt{tslmm}}^2$ is an unbiased estimator of the mutational variance $\sigma_m^2$.
The additive genetic variance $\var{\bg}$ exhibits substantial variability but it is also possible to estimate $\sigma_m^2$ because the average additive variance over replicates is $4N_e \sigma_m^2$.
Based on this observation, one might be tempted to estimate $\var{\bg}$ with $\widehat{\tau}^2_{\texttt{tslmm}}$ through $4N_e \cdot P \widehat{\tau}^2_{\texttt{tslmm}}$.
However, the two quantities are very weakly correlated (\textbf{Figure \ref{fig:mutationalvariance}B}).
Thus, even in the simplest setting of a panmictic population in an equilibrium, 
evolving with constant size, 
in which a closed-form formula,
$\E{\var{\bg}} = 4N_e\sigma_m^2$, between average additive genetic variance $\E{\var{\bg}}$ and $\sigma_m^2 = P \tau^2$ exists, 
one cannot reliably obtain $V_G$ nor $h^2$ out of ARG-LMM.
In more realistic evolutionary models lacking this simple relationship, the attempt to recover $\var{\bg}$ from $\widehat{\tau}^2_{\texttt{tslmm}}$ is even more hopeless.
In sum, the estimates from ARG-LMM based models should be understood as the mutational variance $\sigma_m^2$ divided by the length of the genome and one should refrain from interpreting it as the additive variance.
One can instead estimate the mutational heritability $h^2_m = \sigma_m^2 / \sigma^2$ \citep{lynch1998quant}.

\begin{figure}[ht]
    \centering
    \includegraphics[width=\linewidth]{img/mutationalvariance.pdf}
    \caption{Mutational and additive genetic variances. 10 freely recombining loci were simulated with mutation rate $\mu=10^{-4}$ and a constant population size $N_e=1000$ were simulated 100 times. (A) Additive variance $\var{\bg}$ and \texttt{tslmm} $\widehat{\tau}^2$ can both estimate the mutational variance. (B) The correlation between additive variance and \texttt{tslmm} is weak.}
    \label{fig:mutationalvariance}
\end{figure}




